% !TeX root = ../main.tex
% 3. fejezet — Rendszerarchitektúra
% A ftsrg/thesis-template-latex sablonba \include-dal illesztendő tartalomfájl.
% Hivatkozott források: Backend/HUNting.API/Program.cs, Controllers/*, Services/*,
%   Frontend/HUNting/src/lib/apiClient.ts, services/weatherService.ts, pages/Map/MapPage.tsx,
%   components/Chatbot.tsx, thesis/planning/feature_matrix.md

\chapter{Rendszerarchitektúra}
\label{chap:architektura}

Ez a fejezet madártávlatból mutatja be a HUNting felépítését: a rétegeket és felelősségeiket, a
rétegek közötti kommunikáció szabályait, valamint azokat a tervezési döntéseket, amelyek a
rendszer jelenlegi formáját meghatározzák. A cél a mérnöki szempontból pontos minősítés: a
HUNting \emph{rétegzett monolit} -- nem mikroszolgáltatás-architektúra és nem is a Clean
Architecture értelmében vett, szigorúan elválasztott rétegrendszer --, és a fejezet azt is
megindokolja, hogy a projekt méretéhez és céljaihoz miért ez az ésszerű választás. A rétegek
\emph{belső} felépítését (adatbázis-kontextus, DI, middleware, kliensváz) a
\ref{chap:alapok}.~fejezet részletezi; itt a rétegek \emph{közötti} viszonyokra koncentrálunk.

\section{A háromrétegű architektúra áttekintése}
\label{sec:arch-attekintes}

A rendszer klasszikus háromrétegű felépítést követ, amelyet
a~\ref{fig:architektura}.~ábra szemléltet. A \emph{megjelenítési réteg} egy React-alapú
egyoldalas alkalmazás (single-page application, SPA): a böngésző egyszer tölti le a statikus
HTML-, JavaScript- és CSS-állományokat, ezt követően minden képernyőváltás kliensoldali
útvonalkezeléssel, minden adatművelet aszinkron HTTP-hívással történik. Az \emph{üzleti réteg}
egy ASP.NET Core 8 Web API, amely állapotmentes REST-végpontokon szolgálja ki a klienst: itt
történik a hitelesítés és a jogosultság-ellenőrzés, az üzleti szabályok érvényesítése és a
feltöltött fájlok kezelése. Az \emph{adatréteg} egy Microsoft SQL Server adatbázis (fejlesztői
környezetben LocalDB), amelyet a backend kizárólag az Entity Framework Core objektum-relációs
leképezőn keresztül ér el -- kézzel írt SQL a kódbázisban nem fordul elő.

% TODO: A-01 ábra (figures/a-01-architektura.pdf) — React SPA ↔ REST/JSON ↔ ASP.NET Core API ↔
% EF Core ↔ MSSQL; külső API-k: OpenAI (szerveroldalról), Google Maps + OpenWeather (kliensről).
\begin{figure}[htbp]
	\centering
	\fbox{\parbox[c][7cm][c]{0.9\textwidth}{\centering A-01: A HUNting háromrétegű
	architektúrája a külső szolgáltatásokkal (elkészítendő)}}
	\caption{A rendszer architektúrája: rétegek és külső integrációk}
	\label{fig:architektura}
\end{figure}

A repository ezt a felépítést két fő projektben tükrözi: a \texttt{Backend/HUNting.API} az
API-t és az adatelérést, a \texttt{Frontend/HUNting} a kliensalkalmazást tartalmazza; mindkettő
mellett saját tesztprojekt, illetve tesztkönyvtár áll (\ref{chap:teszteles}.~fejezet). A két
projekt között fordítási idejű függőség nincs: a kapcsolatot kizárólag a HTTP-határfelület és
a JSON-formátum jelenti, a kliensoldali TypeScript-típusok a backend válaszformátumait kézzel
tükrözik.

A képet a külső szolgáltatások teszik teljessé, és itt egy architekturális szempontból fontos
aszimmetria figyelhető meg. Az OpenAI Chat Completions API-t a backend hívja (a chatbot-kérések
a saját \texttt{/api/Chat} végponton keresztül, szerveroldali API-kulccsal jutnak el az OpenAI
felé), a Google Maps JavaScript API-t és az OpenWeather API-t viszont közvetlenül a kliens éri
el, kliensoldali kulcsokkal. Ez utóbbi egyszerűbb megvalósítást adott, ára azonban a kulcsok
böngészőbeli kitettsége; a kompromisszumot a~\ref{chap:chatbot}.~fejezet elemzi.

A rendszer szereplőit és fő használati eseteit a~\ref{fig:usecase}.~ábra foglalja össze. A
\emph{vendég} böngészheti a hirdetéseket, a társaságkatalógust és az időjárásoldalt, valamint
használhatja a chatbotot; a \emph{vadász} (Person fiók) emellett hirdetést ad fel, társasági
tagságot igényel, naplót vezet és térképi objektumokat kezel; a \emph{társasági fiók} (Society)
a tagfelvételt bírálja el, bejegyzéseket és tagokat kezel; az \emph{adminisztrátor} a teljes
felhasználó- és hirdetésállomány felett rendelkezik.

% TODO: U-01 ábra (figures/u-01-usecase.pdf) — Use Case diagram: Vendég, Vadász (Person),
% Társasági fiók (Society), Adminisztrátor aktorok a 6 funkcióterület fő eseteivel.
\begin{figure}[htbp]
	\centering
	\fbox{\parbox[c][7cm][c]{0.9\textwidth}{\centering U-01: A HUNting használati esetei
	(elkészítendő)}}
	\caption{A rendszer szereplői és fő használati esetei}
	\label{fig:usecase}
\end{figure}

\section{REST API és adatfolyamok}
\label{sec:arch-rest}

A kliens és a szerver közötti határfelület a REST elveit követi~\cite{fielding2000}. Az
erőforrásokat főnévi útvonalak azonosítják az egységes \texttt{api/[controller]} sémában
(\texttt{/api/Ads}, \texttt{/api/Societies}, \texttt{/api/HuntingLogs}~stb.), a műveletet pedig
a HTTP-ige hordozza: a \texttt{GET} lekérdez, a \texttt{POST} létrehoz (illetve alárendelt
akciókat indít, például \texttt{POST /api/Societies/\{id\}/join}), a \texttt{PUT} teljes
módosítást végez, a \texttt{DELETE} töröl. A válaszok a szabványos státuszkódokkal élnek:
sikeres lekérdezésre 200, létrehozásra 201 (\texttt{Created}, a létrejött erőforrásra mutató
hivatkozással), tartalom nélküli sikerre 204; a hibás kérés 400-at, a hiányzó vagy érvénytelen
token 401-et, az elégtelen jogosultság 403-at, a nem létező erőforrás 404-et kap. A
reprezentáció mindkét irányban JSON; a szerializáció globális beállításait (karakterlánc-alapú
enumok, körkörös hivatkozások elvágása, \texttt{null} mezők elhagyása) a
\ref{chap:alapok}.~fejezet mutatta be.

Egy tipikus kérés életútját a~\ref{fig:request-flow}.~ábra követi végig. A kliensoldali
funkciómodul az \texttt{apiClient} példányon indítja a hívást, amelynek kérés-elfogója befűzi
a bearer-tokent. A kérés a szerveren a middleware-csővezetéken halad át: a CORS-ellenőrzés után
a hitelesítési middleware validálja a tokent és felépíti a kérés identitását, a
jogosultságkezelés kiértékeli az \texttt{[Authorize]} szabályokat, majd az útvonalillesztés a
megfelelő kontrollerakciót hívja. Az akció az erőforrás-szintű ellenőrzések után a
szolgáltatásrétegen keresztül az EF Core-ral fordul az adatbázishoz, az eredmény pedig
JSON-válaszként utazik vissza, ahol a React Query gyorsítótárba kerül, és a komponensek
deklaratívan megjelenítik. Hiba esetén a kliens válasz-elfogója lép közbe: a 401-es választ
globális kijelentkeztetésként kezeli, minden mást egységes hibaüzenetté normalizál.

% TODO: A-02 ábra (figures/a-02-keres-eletut.pdf) — folyamatábra: apiClient interceptor → CORS →
% authentikáció → autorizáció → controller → service → EF Core → MSSQL → JSON válasz → React Query.
\begin{figure}[htbp]
	\centering
	\fbox{\parbox[c][6.5cm][c]{0.9\textwidth}{\centering A-02: Egy REST-kérés életútja a
	kliens-elfogótól az adatbázisig és vissza (elkészítendő)}}
	\caption{Egy hitelesített REST-kérés feldolgozási folyamata}
	\label{fig:request-flow}
\end{figure}

Az adatfolyamok szempontjából a rendszer három forgalomtípust bonyolít. A \emph{strukturált
adatforgalom} a fenti JSON-séma szerint zajlik. A \emph{bináris forgalom} a képfeltöltésekre és
-letöltésekre korlátozódik: a feltöltés \texttt{multipart/form-data} formátumú POST-kérés, a
letöltés pedig a statikus fájl-middleware-en keresztül, a JSON-válaszokban kapott URL-ek alapján
történik. Végül a \emph{külső forgalom} a chatbot-kérések továbbítása az OpenAI felé
(szerveroldalon), illetve a térkép- és időjárásadatok lekérése (kliensoldalon).

\section{Tervezési döntések és kompromisszumok}
\label{sec:arch-dontesek}

\paragraph{Monolit a mikroszolgáltatások helyett.} A backend egyetlen telepítési egység, közös
adatbázissal. Egy hat funkcióterületből álló, egy fejlesztő által készített alkalmazásnál a
mikroszolgáltatásokra bontás (szolgáltatásonkénti adatbázis, szolgáltatásközi kommunikáció,
elosztott tranzakciók) aránytalan üzemeltetési és fejlesztési többletköltséget jelentene,
miközben az előnyei -- független skálázás és telepítés -- itt nem jelentkeznének. A monoliton
belül a funkcióterületek a kontroller--szolgáltatás--modell hármasok mentén így is jól
elkülönülnek, ami a későbbi kiemelést nem zárja ki.

\paragraph{Rétegzett monolit, nem Clean Architecture.} A rétegek elválasztása pragmatikus, nem
dogmatikus: az entitásosztályok egyszerre szolgálnak adatbázis-leképezésként és -- több végpontnál
-- API-válaszformátumként, a szolgáltatásréteg pedig nem mindenhol van jelen (a hírek kontrollere
közvetlenül a kontextust használja). Külön tartományi réteg, interfészprojekt vagy
DTO-leképező keretrendszer (például AutoMapper) nincs: ahol a válasz formátuma eltér az
entitástól (például a hitelesítési válaszok), ott kézi leképezés történik. Ez a döntés
csökkenti az absztrakciós rétegek számát és a kód mennyiségét, cserébe a séma és az API
összefonódik: egy entitásmódosítás közvetlenül kihat a válaszformátumra. A dolgozat ezt az
őszinte minősítést következetesen tartja: ahol egy fejezet a rétegzéstől való eltérést talál,
ott azt jelzi (így a részleges repository-mintát a~\ref{sec:alapok-backend}.~alfejezet
tárgyalja~\cite{fowler2002}).

\paragraph{Oldalszintű védelem a kliensen.} A kliensoldali hozzáférés-védelem nem központi
útvonalőrrel, hanem az egyes oldalkomponensekben valósul meg
(\ref{sec:alapok-frontend}.~alfejezet). Ez a megközelítés egyszerű és oldalanként testre
szabható (eltérő átirányítási célok, szerepkörfüggő nézetek), hátránya a szabályok szétszórtsága.
Hangsúlyozandó ugyanakkor az architektúra alapelve: a kliensoldali védelem kizárólag felhasználói
élményt szolgál, a tényleges hozzáférés-szabályozás minden esetben a szerveren történik -- a
kliens megkerülésével indított kérések is a JWT-validáción és a szerepkör-, illetve
tulajdonos-ellenőrzéseken buknak el.

\paragraph{A funkcióközpontú dokumentálás indoklása.} A dolgozat szerkezete maga is tervezési
döntés. Mivel a rendszer értéke a funkcióiban rejlik, és minden funkció ugyanazon a három
rétegen ível át, a dokumentáció nem rétegenként (külön adatbázis-, backend- és
frontend-fejezetekben), hanem funkciónként halad: az 5--10.~fejezet mindegyike az adatmodelltől
a felhasználói felületig mutat be egy-egy funkciót. Így az összetartozó döntések egy helyen
olvashatók, az ismétlődő infrastruktúra pedig a \ref{chap:alapok}.~fejezetben, egyszer szerepel.
